\section{Discussion}
\label{sec:discussion}

\subsection{What the inflight uplift really means}

The most striking single number in the paper is the LightGBM scenario-A$\to$B PR-AUC jump from~\num{0.520} to~\num{0.951}. It's worth being precise about what this means. The lag-feature ablation (\cref{tab:lag-ablation}) shows that \emph{both} scenarios collapse to PR-AUC~\num{0.355} when the lag and inflight features are removed -- a static-features-only model. Since A retains the lag features but not the inflight features, and B retains both, the ablation says:

\begin{itemize}
    \item Lag features alone (no inflight) $\Rightarrow$ \num{0.520} PR-AUC (scenario A)
    \item Lag + 3 inflight features $\Rightarrow$ \num{0.952} PR-AUC (scenario B)
    \item Neither $\Rightarrow$ \num{0.355} PR-AUC (the static-features baseline)
\end{itemize}

So the marginal value of the three inflight features (\texttt{prev\_stop\_actual\_delay}, \texttt{cum\_actual\_delay\_so\_far}, \texttt{max\_actual\_delay\_so\_far}) is +0.43 PR-AUC -- an enormous lift, dominated by \texttt{prev\_stop\_actual\_delay} per the SHAP analysis (\cref{fig:global-bar}b). Operationally this confirms that \emph{within-run} delay propagation is highly Markovian: knowing the prior stop's actual delay tells you almost everything about whether the next stop will be late.

\subsection{The \emph{within-NL test} vs \emph{cross-country transfer} gap}

The cause-of-disruption v2 model achieves F1=\num{0.509} on the rare \emph{weather} class within the Netherlands -- a non-trivial result. The same model predicts weather for only~0.1\,\% of Finnish disrupted stops despite Finland having an extreme winter weather regime. This is not a model-quality issue; it's a domain-shift issue. The features that the NL model learned to associate with weather (Open-Meteo daily precipitation aggregates around a calibrated threshold, snow depth in centimetres on Dutch winters that cap around \SI{5}{cm}) do not match Finland's FMI hourly observations on a different numerical scale and Finland's snow depth peaks of \SI{40}{cm}+. The v2 transfer to Finland is much more domain-plausible than v1 (no more rolling-stock + logistical pathology), but the rare-but-important weather class is still missed because the features themselves don't transfer.

\subsection{The role of the Knowledge Graph}

The KG is not a model in its own right; it's an integration layer between the SHAP report and the spatial-temporal context. After the SHAP top-driver analysis flags \texttt{prev\_stop\_actual\_delay} as the dominant inflight feature, the bridge query (\cref{lst:kg-bridge}) returns the surrounding station's degree, recent fault count, and adjacency neighbours -- giving the operations analyst a structured view of \emph{why} the model is confident. The fact that the bridge query runs in \SI{52}{ms} on \num{16.59}~M relationships means it's responsive enough to drop into an interactive disruption-monitoring dashboard.

\subsection{Limitations}

\begin{enumerate}
    \item \textbf{BiLSTM was deferred} (\cref{sec:models}). The architecture is sound at synthetic scale but \texttt{\_pack\_sequences} OOMs on full data. A batched generator-based packer would solve this.
    \item \textbf{Italy zero-fill weather bug.} \texttt{lib\_italy.py} zero-fills missing weather columns instead of preserving \texttt{NaN}, which contaminates downstream feature distributions and is the root cause of the v2 IT \emph{engineering work} mode collapse.
    \item \textbf{Sparksee runs in evaluation mode.} The Sparksee 5.2.3 jar is pulled from Maven Central and runs without a licence string, which caps each edge type at \SI{1}{M} entries. The build script auto-falls back to the FI\_HKH hub subgraph (98.5~K \texttt{STOPS\_AT} edges); holders of an academic / commercial licence can drop a \texttt{sparksee.cfg} with their licence string and load the full~\SI{16.59}{M}-edge graph without any code changes.
    \item \textbf{Cause label noise.} Only~32\,\% of NL disrupted stops have a matching RDT incident-log record. The remaining stops are disrupted but unlabelled -- our ``labelled set'' is a biased subsample of the full disrupted population.
    \item \textbf{Static cross-country feature scales.} Open-Meteo (daily aggregates) and FMI (hourly observations) produce numerically different precipitation distributions, breaking transfer of weather-related cause predictions.
\end{enumerate}

\subsection{Reproducibility}

The whole pipeline is reproducible from raw open data with three shell commands; total wall-clock time on a 4-core CPU + Tesla T4 box is~$\sim$2.5 hours (Phase 1: 35\,min; Phase 2: 25\,min; Phase 3: 35\,min; Phase 4: 35\,min; KG: 10\,min; cause v2: 20\,min). The 32-test pytest suite runs in~25\,s. All data sources are public or publicly-licensable (Trenitalia~\cite{trenitalia2025figshare}~CC-BY 4.0; FI-TW~\cite{borin2025fitw}~GPL-3 + Kaggle account; RDT~\cite{rdt2024opendata}~CC-BY 4.0; Open-Meteo~\cite{zippenfenig2023openmeteo}~free).
